% 1. はじめに: 問い、単位、重力の微細構造定数
\section{はじめに: 何を調べるのか}\label{sec:intro}

\subsection{研究の問い}

質量 $\mu\sim10^{-12}\eV$ の非常に軽いスカラー粒子(ALP、axion-like particle)があると、回転するブラックホール(BH)のまわりに、その粒子が非常に多く集まった「雲」ができる。雲は BH の回転のエネルギーと角運動量を吸い取って育つ。そのため、この粒子があれば、ある範囲の質量の BH は速く回っていられない。

重力波の観測(LIGO--Virgo--KAGRA、カタログ GWTC-5)では、合体する BH の質量とスピンが測られている。そこから、この粒子について逆の結論が出ている。
\begin{itemize}
\item 速く回る BH がいるので、この質量の粒子は排除される。
\item BH のスピンが質量ごとにある値に集まっているように見えるので、この質量の粒子がいる証拠がある。
\end{itemize}
研究の問いは、粒子どうしの相互作用(自己相互作用)を入れるとこの二つが両立するか、である。自己相互作用があると、雲が大きく育つ前に飽和し、BH が回転を失う速さが遅くなる。若い BH は回転を保ち、年をとった BH だけが回転を失う、ということが起こりうる。

このノートでは、この問いに答えるのに要る物理を順に導く。
\begin{enumerate}
\item 回転する BH の性質(\S\ref{sec:kerr})
\item BH から回転を取り出す「超放射」の条件(\S\ref{sec:superradiance})
\item BH のまわりの粒子の束縛状態「重力原子」と、その成長率(\S\ref{sec:atom})
\item 雲の成長と BH の回転の減少(\S\ref{sec:cloud})
\item 自己相互作用と、それが起こす準位間の遷移(\S\ref{sec:selfint})
\item 準平衡と、回転を失う速さ(\S\ref{sec:equilibrium})
\item 吸収で回転が戻らないこと(\S\ref{sec:norestore})
\item 合体で生まれる 2 世代目の BH のスピン(\S\ref{sec:secondgen})
\item これらを組み合わせた「窓」(\S\ref{sec:window})
\end{enumerate}

\begin{figure}[H]
\centering
% BH と雲の大きさ(模式図)
\begin{tikzpicture}[scale=1.0]
  % 雲(同心円のぼかし)
  \foreach \r/\o in {3.0/0.05, 2.6/0.08, 2.2/0.11, 1.8/0.14, 1.4/0.12, 1.0/0.08}{
    \fill[cA, opacity=\o] (0,0) ellipse ({\r} and {0.55*\r});
  }
  % BH
  \fill[ink] (0,0) circle (0.12);
  \draw[ink2, thin] (0,0) circle (0.17);
  % 回転の矢印
  \draw[-{Stealth[length=2mm]}, ink2] (0.32,0.1) arc (10:300:0.32 and 0.18);
  % 寸法
  \draw[{Stealth[length=2mm]}-{Stealth[length=2mm]}, ink2] (0,-0.25) -- (2.2,-1.21) node[midway, below right, font=\small, color=ink] {$r_c\simeq r_g/\alpha^2$};
  \node[font=\small, anchor=west] at (0.25,0.45) {BH($r_g$)};
  \node[font=\small, color=cA!80!black] at (2.9,1.45) {粒子の雲};
  \node[font=\scriptsize, color=ink2, align=left, anchor=west] at (3.4,0.2) {$\alpha=0.1$ なら\\$r_c\approx100\,r_g$};
\end{tikzpicture}

\caption{BH と雲の大きさの関係(模式図)。雲の大きさ $r_c\simeq r_g/\alpha^2$ は、$\alpha\sim0.1$ では BH の大きさ $r_g$ の 100 倍ほどになる。雲のごく一部が地平線にかかり、そこで BH と角運動量をやりとりする。}
\label{fig:bhcloud}
\end{figure}

\subsection{単位と記号}

\paragraph{自然単位}
このノートでは、ことわりがない限り $\hbar=c=1$ とする。このとき質量・エネルギー・振動数・長さの逆数はすべて同じ次元をもつ。数値を出すときは次の換算を使う。
\begin{equation}
\hbar=6.582\times10^{-16}\eV\,\mathrm{s},\qquad
\frac{\mu c^2}{\hbar}=1.519\times10^{3}\,\mathrm{s}^{-1}\times\frac{\mu}{10^{-12}\eV}.
\label{eq:hbar}
\end{equation}

\paragraph{BH の尺度}
質量 $M$ の BH は、長さ $r_g$ と時間 $t_g$ の尺度をもつ。
\begin{equation}
r_g\equiv\frac{GM}{c^2}=1.477\,\mathrm{km}\times\frac{M}{M_\odot},\qquad
t_g\equiv\frac{GM}{c^3}=4.925\,\mu\mathrm{s}\times\frac{M}{M_\odot}.
\label{eq:rg}
\end{equation}
$r_g$ を「重力半径」と呼ぶ(Schwarzschild 半径 $2r_g$ の半分)。BH の章(\S\ref{sec:kerr})では、さらに $G=1$ とおき、$M=1$ を長さの単位にすることがある。

\subsection{重力の微細構造定数 $\alpha$}\label{sec:alpha}

この研究のすべての式に現れる無次元量が
\begin{equation}
\alpha\equiv\frac{GM\mu}{\hbar c}
\label{eq:alpha}
\end{equation}
である。意味は三通りに読める。

\paragraph{読み方 1: 二つの長さの比}
粒子の換算 Compton 波長は $\lambda_c=\hbar/(\mu c)$ である。これと重力半径の比をとると
\begin{equation}
\frac{r_g}{\lambda_c}=\frac{GM/c^2}{\hbar/(\mu c)}=\frac{GM\mu}{\hbar c}=\alpha.
\end{equation}
$\alpha\ll1$ なら、粒子の波は BH よりずっと大きい。

\paragraph{読み方 2: 時間の比}
式 \eqref{eq:alpha} は、粒子の振動数 $\mu c^2/\hbar$ と BH の時間 $t_g$ の積でもある。
\begin{equation}
\alpha=\frac{\mu c^2}{\hbar}\times\frac{GM}{c^3}=\frac{\mu c^2}{\hbar}\,t_g.
\label{eq:alpha_time}
\end{equation}
数値は式 \eqref{eq:hbar}、\eqref{eq:rg} から
\begin{equation}
\alpha=1.519\times10^{3}\,\mathrm{s}^{-1}\times4.925\times10^{-6}\,\mathrm{s}\times\frac{M}{M_\odot}\frac{\mu}{10^{-12}\eV}
=0.0748\times\frac{M}{10\,M_\odot}\frac{\mu}{10^{-12}\eV}.
\label{eq:alpha_num}
\end{equation}

\paragraph{読み方 3: 水素原子の $\alpha_\mathrm{em}$ の重力版}
水素原子では、電子(質量 $m_e$)が Coulomb ポテンシャル $V=-e^2/(4\pi\varepsilon_0r)\equiv-\alpha_\mathrm{em}\hbar c/r$ に束縛される。BH のまわりの粒子(質量 $\mu$)は、遠方では Newton のポテンシャル
\begin{equation}
V(r)=-\frac{GM\mu}{r}=-\alpha\frac{\hbar c}{r}
\label{eq:Vgrav}
\end{equation}
に束縛される。$\alpha_\mathrm{em}$ のところに $\alpha$ が入るだけなので、\S\ref{sec:atom} で水素原子の結果をそのまま使える。

\begin{numcheck}
GW241011 の主星($M=19.6\Msun$)と $\mu=0.71\times10^{-12}\eV$ では、式 \eqref{eq:alpha_num} から
\[
\alpha=0.0748\times1.96\times0.71=0.104.
\]
研究ノートの値(E2、E3)と一致する。
\end{numcheck}

\begin{tocode}
コード: \texttt{research/calc/sr.py} の \texttt{alpha\_of(M, mu\_eV)} が式 \eqref{eq:alpha_time} を計算する。定数は \texttt{TSUN} $=GM_\odot/c^3$、\texttt{HBAR\_EVS} $=\hbar$ [eV s]。
\end{tocode}

\begin{summary}
\begin{itemize}
\item 長さの尺度は $r_g=GM/c^2$、時間の尺度は $t_g=GM/c^3$。
\item $\alpha=GM\mu/(\hbar c)=0.0748\,(M/10M_\odot)(\mu/10^{-12}\eV)$ は、BH の大きさと粒子の波長の比であり、重力原子の微細構造定数である。
\end{itemize}
\end{summary}
